% Laget automatisk av elevark/lag_elevark.py fra opplegg/sentripetalkraft-snor.html. Endringer her blir overskrevet.
\documentclass[11pt]{article}
\usepackage{elevark}
\begin{document}
\elevarkhode{Fysikk 2 · Sirkelbevegelse}{Sentripetalkraft med snor og rør}{https://fysikklab.no/opplegg/sentripetalkraft-snor.html}{Opplegget på nett}

En gummipropp er bundet i en snor som går gjennom et rør. I den andre enden henger et lodd. Sving proppen rundt over hodet slik at loddet holder seg i ro. Da er det tyngden av loddet som holder proppen i sirkelbanen. Stemmer det med $F = mω^{2}L$?

\section*{Hensikt}

\begin{itemize}
\item Undersøke sammenhengen mellom kraften og vinkelfarten i en sirkelbevegelse.
\item Sammenligne snordraget med sentripetalkraften $mω^{2}L$.
\end{itemize}

\section*{Utstyr}

\begin{itemize}
\item En gummipropp, omtrent 1,5~m tynn, sterk snor og et glatt rør (for eksempel en plastpenn uten innmat eller et glassrør med tape rundt).
\item Loddholder og lodd (50–250~g), en binders eller tape som merke på snora, stoppeklokke og vekt.
\item Vernebriller. Hold god avstand til andre, og sving over hodehøyde.
\end{itemize}

\section*{Teorien}

Proppen går i en vannrett sirkel, men snora heller litt nedover fra røret til proppen. Kall snorlengden fra røret til proppen $L$, og vinkelen under vannrett $α$. Radien er da $r = L \cos  α$. Den vannrette komponenten av snordraget $S$ gir sentripetalkraften: $S \cos  α = mω^{2}r = mω^{2}L \cos  α$. Vinkelen forsvinner, og $S = mω^{2}L$. Når loddet henger i ro, er $S = Mg$. Altså skal $Mg = mω^{2}L$.

\section*{Framgangsmåte}

\begin{enumerate}
\item Vei proppen. Tre snora gjennom røret og fest proppen i den ene enden og loddholderen i den andre.
\item Velg $L$, for eksempel 60~cm fra toppen av røret til midten av proppen. Sett en binders på snora rett under røret som merke.
\item Hold røret loddrett og sving proppen rundt. Juster farten til bindersen holder seg rett under røret, slik at $L$ holder seg fast.
\item Ta tiden for 20 runder når farten er jevn.
\item Gjenta med fem ulike lodd, og samme $L$.
\end{enumerate}

\section*{Oppgaver}

\begin{oppgaveliste}
\oppgave \fase{gjett}{Gjett}Dere firedobler loddet og holder $L$ fast. Hvor mye fortere må proppen gå rundt?
\svar
\oppgave \fase{test}{Mål}Gjør målingene, eller trykk på «Fyll inn eksempel». Ligger punktene nær linjen med stigningstall 1?
\svar
\oppgave \fase{forklar}{Forklar}Vis at snordraget er $S = mω^{2}L$, uansett hvor mye snora heller. Hvorfor kan snora aldri være helt vannrett?
\svar
\oppgave \fase{test}{Regn}Med $m = 25\,\mathrm{g}$ og $L = 60\,\mathrm{cm}$: Hvor lang tid bruker proppen på én runde når loddet er 100~g? Hvor stor er farten til proppen?
\svar
\oppgave \fase{forklar}{Vurder}Hvorfor blir $mω^{2}L$ ofte litt mindre enn Mg? Tenk på røret, snora og hvordan dere måler $L$.
\svar
\oppgave \fase{utfordring}{Utfordring}Hvor stor vinkel $α$ heller snora med når loddet er 100~g? Bruk den loddrette komponenten av snordraget.
\svar
\end{oppgaveliste}

\end{document}
